\documentclass[a4paper]{article}

\usepackage[english]{babel}
%\usepackage[utf8x]{inputenc}
\usepackage{amsmath}
\usepackage{graphicx}
\usepackage[colorinlistoftodos]{todonotes}
\usepackage[all]{xy}
\usepackage{dsfont}
\usepackage{amssymb}
\usepackage{bbold}
\usepackage{amsthm}

\usepackage{tikz}
\usepackage{verbatim}
\usetikzlibrary{matrix}
\usetikzlibrary{cd}

\usepackage[autostyle=true]{csquotes}
\usepackage{pdfpages}
\usepackage{hyperref}
\usepackage{fancyhdr}



\newtheorem{theorem}{Theorem}[section]
\newtheorem{proposition}{Proposition}[theorem]
\newtheorem{lemma}{Lemma}[theorem]
\newtheorem{corollary}{Corollary}[theorem]
\newtheorem*{definition}{Definition}


\newcommand{\norm}[1]{\left\lVert#1\right\rVert}

\newcommand{\Rn}{\mathbb{R}^n}


\newcommand{\too}[1]{\overset{#1}{\longrightarrow}}
\newcommand{\hookar}[1]{\overset{#1}{\hookrightarrow}}
\newcommand{\colim}[1]{ \underset{#1}{\texttt{colim}} \text{ } }
\newcommand{\fset}[1]{ \{ #1_0 , \dots , #1_n \} }
\newcommand{\ddt}{\frac{d}{dt}}


\newcommand{\N}{\mathbb{N}}
\newcommand{\Z}{\mathbb{Z}}







\fancyhead[L]{\includegraphics[scale=0.5]{UFT.png}}

\begin{document}

\begin{center}
    \Large{\textbf{Propuestas}}
\end{center}

\setlength{\parindent}{0cm}


\textbf{Ejercicio 1:} Responda las siguientes preguntas \\


a) Determine la ecuación general del plano que pasa por el punto $(1,2,0)$ y que contiene a la recta de ecuación vectorial $L_1: (2,-1,3)+t(1,0,-2)$. \\


b) Demuestre que el plano  encontrado en a) es paralelo a la recta $L_2: (1,2,0)+t(3,3,-11)$ \\ 

c) Demuestre que el plano  encontrado en a) es perpendicular a la recta $L_3$ de ecuación general 

\begin{align*}
    \frac{x+1}{3} = \frac{2y-3}{5} =\frac{2z+3}{3}
\end{align*}











\textbf{Ejercicio 2:} Calcule los siguientes límites, en caso de no existir demuéstrelo \\

a) $\displaystyle \lim_{(x,y) \to (0,0)} \frac{x^2y}{\sqrt{x^2+y^2}}$ \\


b) $\displaystyle \lim_{(x,y) \to (3,3)} \frac{9x^2-y^4}{3x-y^2}$ \\

c) $\displaystyle \lim_{(x,y) \to (0,0)} \frac{x^3-3x^2y}{(x^2+y^2)^{3/2}}$ \\ \\ \\


\textbf{Ejercicio 3:} Considere el cilindro $x^2+y^2=9$ y el $x-y+z=4$, encuentre una parametrización $r(t)$ de la intersección de estas dos superfices, luego demuestre que $r(0)$ es perpendicular a la recta $L$ de ecuación

\begin{align*}
    \frac{x-1}{2} = \frac{1-y}{5} = \frac{2-z}{6}
\end{align*}





\end{document}